% SOURCE_STATEMENT: "On every compact negatively curved manifold, do all high-energy Laplace eigenfunctions equidistribute?"
% Preparation and final self-audit: CONTRIBUTING.md prompts C and D; source/audit artifacts retained in the local physics_statements_2026_10_11 batch.
% !TeX program = lualatex
\documentclass[11pt]{article}
\usepackage[margin=25mm]{geometry}
\usepackage{fontspec}
\setmainfont{Latin Modern Roman}
\usepackage{amsmath,amssymb,mathtools}
\usepackage{unicode-math}
\setmathfont{Latin Modern Math}
\usepackage{xurl,hyperref}
\hypersetup{colorlinks=true,urlcolor=blue}
\setcounter{secnumdepth}{0}
\setlength{\parindent}{0pt}
\setlength{\parskip}{5pt}
\setlength{\emergencystretch}{3em}
\newcommand{\sref}[2]{\hyperlink{source:#1:#2}{[S#2]}}
\newcommand{\cataloguescope}[1]{\paragraph{Recorded target.}#1}
\newcommand{\R}{\mathbb R}
\newcommand{\C}{\mathbb C}
\newcommand{\N}{\mathbb N}
\newcommand{\Z}{\mathbb Z}
\newcommand{\Q}{\mathbb Q}
\newcommand{\F}{\mathbb F}
\newcommand{\eps}{\varepsilon}
\begin{document}
\section{The Quantum Unique Ergodicity Conjecture}\label{statement}

\subsection{Definitions and mathematical statement}

\paragraph{Shared notation.}$\mathbb N=\{1,2,\ldots\}$; the standard number fields have their usual meanings. All additional conventions are stated locally.


\paragraph{Definitions and conventions.}
Let $(M,g)$ be a compact boundaryless negatively curved Riemannian manifold and $-\Delta_g\varphi_j=\lambda_j^2\varphi_j$, with $\|\varphi_j\|_{L^2}=1$ and $\lambda_j\to\infty$. Configuration-space equidistribution means $\int_M a|\varphi_j|^2\,d\mathrm{vol}_g\to \mathrm{vol}_g(M)^{-1}\int_M a\,d\mathrm{vol}_g$ for every continuous test function $a$. The standard stronger microlocal QUE formulation concerns convergence of semiclassical measures on the unit cotangent bundle to normalized Liouville measure.

\paragraph{Target retained from the canonical source catalogue.}
\paragraph{Catalogue wording retained for traceability.}
On every compact negatively curved manifold, do all high-energy Laplace eigenfunctions equidistribute?

\paragraph{Source scope and qualification.}
The supplied sentence asks equidistribution of all eigenfunctions, without discarding an exceptional subsequence. The displayed spatial criterion explains that wording; the microlocal version is identified separately and not silently treated as the same weaker target.
\subsection{Short English statement}
Must every high-energy eigenfunction sequence on a compact negatively curved manifold become uniformly distributed?
\subsection{Sources}
\begin{itemize}
\item[\hypertarget{source:statement:1}{S1}] ULAM.AI and the UnsolvedMath Contributors, UnsolvedMath, catalogue entry. Catalogue curation under CC BY 4.0; definitions and English exposition added locally.
\url{https://huggingface.co/datasets/ulamai/UnsolvedMath}.
\item[\hypertarget{source:statement:2}{S2}] Free Probability Theory (2018), Oberwolfach Reports, DOI 10.4171/owr/2018/53, pp. 3147--3215. Primary report accessed 2026-10-11; relevant PDF pages 55.
\url{https://doi.org/10.4171/owr/2018/53}.
\item[\hypertarget{source:statement:3}{S3}] Lindenstrauss (2006) (catalogue-listed context; not independently consulted in this preparation).
\url{https://annals.math.princeton.edu/2006/163-1/p03}.
\item[\hypertarget{source:statement:4}{S4}] Soundararajan (2010) (catalogue-listed context; not independently consulted in this preparation).
\url{https://arxiv.org/abs/0901.4060}.
\end{itemize}
\label{end:statement}
\end{document}
